% ****** Start of file final_project_report.tex ******
%
%   This file follows the same REVTeX style as the user's original
%   Phys765_final_project.zip template.
%

\documentclass[%
 reprint,
 %superscriptaddress,
 %groupedaddress,
 %unsortedaddress,
 %runinaddress,
 %frontmatterverbose,
 %preprint,
 %preprintnumbers,
 %nofootinbib,
 %nobibnotes,
 %bibnotes,
 amsmath,amssymb,
 aps,
 %pra,
 %prb,
 %rmp,
 %prstab,
 %prstper,
 floatfix,
]{revtex4-2}

\usepackage{booktabs}
\usepackage{makecell}
\usepackage{braket}
\usepackage{graphicx}
\usepackage{quantikz}
\usepackage{dcolumn}
\usepackage{bm}
%\usepackage{hyperref}
%\usepackage[mathlines]{lineno}
%\linenumbers\relax

\begin{document}

\preprint{APS/123-QED}

\title{Comparison of Quantum Krylov and UCCSD for Molecular Ground-State Approximation via Subspace Projection}

\author{Jihyeon Park}
\email{park687@wisc.edu}

%\date{\today}

\begin{abstract}
This work compares Quantum Krylov and UCCSD as non-variational basis generators for
molecular ground-state approximation in a projected-subspace framework. Rather
than optimizing one ansatz, I compare how well each basis family generates a
useful low-energy subspace. Across the benchmark set, Quantum Krylov reaches chemical accuracy with
substantially fewer basis states than the non-variational UCC constructions,
showing a stronger expressivity-per-basis trend as system size grows. At the
same time, per-state resource estimates show that Quantum Krylov states are much more
expensive to prepare because they rely on Trotterized Hamiltonian evolution,
whereas selected UCCSD states remain shallow and inexpensive. The resulting
full-stack conclusion is a clear tradeoff: Quantum Krylov is favored at the algorithmic
layer by its smaller required subspace, while selected UCCSD is favored at the
logical and hardware layers by its lower circuit cost.
\end{abstract}

\maketitle

\section{Introduction and Motivation}

Quantum algorithms are widely expected to offer an eventual advantage over classical methods for describing quantum many-body systems, especially in electronic-structure problems where the Hilbert space grows exponentially with system size. Quantum phase estimation (QPE) is one of the representative examples because, in principle, it can recover eigenvalues to high precision when provided with a sufficiently good initial state \cite{abrams1999qpe}. In practice, however, QPE requires deep coherent time evolution and fault-tolerant resources that remain beyond the reach of current hardware for most chemically relevant systems.

These limitations motivated the development of hybrid algorithms such as the variational quantum eigensolver (VQE) \cite{peruzzo2014vqe}. VQE replaces long coherent evolution with repeated preparation and measurement of a parameterized trial state, while a classical optimizer updates the circuit parameters to minimize the measured energy. This hybrid structure is much more compatible with near-term hardware, but it introduces a different bottleneck: the variational optimization itself can be costly, sensitive to initialization, and difficult to scale when the ansatz contains many parameters.

\subsection{Why It Matters}

As an alternative, subspace methods have also been developed for molecular ground-state estimation. Instead of optimizing a single wavefunction, these methods prepare a small set of basis states and solve a projected eigenvalue problem within the span of those states. In this setting, the quality of the chosen basis becomes the central issue. A better basis can improve the accuracy of the projected ground-state estimate, while a lower-cost basis can reduce circuit depth, gate count, and measurement overhead.

This motivates the present study. Among the many basis-construction strategies proposed in the literature, Krylov states generated by Hamiltonian evolution \cite{cortes2022qksd} and UCCSD states generated by chemically motivated single and double excitations \cite{anand2022ucc} are two representative state-of-the-art directions. Krylov uses the Hamiltonian itself to generate candidate directions in Hilbert space, whereas UCCSD uses an excitation structure inherited from coupled-cluster theory. Although many variations of both approaches exist, these two families serve here as representative Hamiltonian-guided and chemistry-guided basis constructions.

\subsection{Problem Statement}

There are several distinct ways to evaluate a quantum state-preparation strategy for ground-state estimation. In a QPE-type setting, a prepared state is used as the initial state for phase estimation, so its quality is judged mainly by its overlap with the target eigenstate. In a VQE-type setting, the state or ansatz is used as the trial form in a variational optimization, so its quality is judged by the minimum energy reachable after parameter updates. In a subspace method, by contrast, a family of prepared states is treated as a basis set, and the quality of that family is judged by the projected eigensolver built from their span.

This work adopts the third viewpoint. The main object of interest in this paper is not just one trial state, but the basis generator itself. Evaluating Krylov and UCCSD through the subspace they generate makes it possible to compare their collective expressivity, which can be richer than what is visible from a single state alone. This choice is especially important for Krylov states: a single Krylov state of the form
\begin{equation}
\ket{\phi(t)} = e^{-it\hat{H}}\ket{\phi_{\mathrm{HF}}}
\end{equation}
has a constant expectation value of $\hat{H}$ under the same Hamiltonian, so it is not meaningful to assess Krylov as a one-state variational ansatz. Krylov is naturally a set-based construction, and subspace projection is therefore the appropriate framework for comparing it with UCCSD as a basis-state generator. The concrete problem studied in this paper is molecular ground-state energy estimation. Starting from the same Hartree-Fock reference and using the same target subspace size $K$, this study asks which basis family yields the lower projected ground-state energy for a given state-preparation cost, with final projected energies compared against exact full configuration interaction (FCI) reference data on benchmark molecules.

\section{Methods / Model Assumptions}

\subsection{Quantum Basis}

\subsubsection{Quantum Krylov Basis}

Quantum Krylov methods construct a basis from repeated Hamiltonian action on a reference state and are naturally suited to projected eigensolvers \cite{cortes2022qksd}. Formally, the Quantum Krylov subspace of order $K$ generated from $\ket{\phi_{\mathrm{HF}}}$ is
\begin{equation}
\mathcal{K}_K(\hat{H}, \ket{\phi_{\mathrm{HF}}}) =
\mathrm{span}\{\ket{\phi_{\mathrm{HF}}}, \hat{H}\ket{\phi_{\mathrm{HF}}}, \hat{H}^2\ket{\phi_{\mathrm{HF}}}, \ldots \}.
\end{equation}
In the circuit-based implementation used here, a time-evolved version is more convenient, and each additional Quantum Krylov basis state is generated from the Hartree-Fock reference using Trotterized real-time evolution,
\begin{equation}
\ket{\phi(t_k)} \approx e^{-it_k\hat{H}}\ket{\phi_{\mathrm{HF}}}.
\end{equation}
The main characteristic of the Quantum Krylov basis is that it is Hamiltonian-guided: the Hamiltonian itself determines the directions added to the span. Its most important continuous parameter is the evolution time $t_k$ assigned to each basis state, while the Trotter order and number of Trotter steps are fixed globally. In principle, the time steps can be selected by grid search or by heuristic spacing rules for time-propagated basis states \cite{stair2020mrsqk,cohn2021qfd}. In this paper, the Quantum Krylov times are chosen as an equally spaced sequence $t_k = k\Delta t$, and the spacing $\Delta t$ is selected by grid search over candidate values. Candidate spacings that lead to a poorly conditioned overlap matrix are rejected, and the final choice is the spacing that gives the lowest projected ground-state energy among the remaining candidates. Quantum Krylov is therefore treated as a non-variational basis generator whose parameter-selection step is restricted to this explicit grid-search protocol rather than to a full continuous optimization loop.

\subsubsection{UCCSD Basis}

UCCSD constructs basis states from a chemically motivated excitation manifold built on top of the Hartree-Fock reference \cite{anand2022ucc}. Its general form is
\begin{equation}
\ket{\phi(\boldsymbol{\theta})}
= e^{\hat{T}(\boldsymbol{\theta})-\hat{T}^{\dagger}(\boldsymbol{\theta})}
\ket{\phi_{\mathrm{HF}}},
\end{equation}
where the cluster operator is truncated to single and double excitations,
\begin{equation}
\hat{T}(\boldsymbol{\theta})
= \sum_{ia} \theta_i^a \hat{a}_a^\dagger \hat{a}_i
+ \frac{1}{4}\sum_{ijab}\theta_{ij}^{ab}
\hat{a}_a^\dagger \hat{a}_b^\dagger \hat{a}_j \hat{a}_i .
\end{equation}
Here $i,j$ denote occupied spin orbitals and $a,b$ denote virtual spin orbitals. The main characteristic of UCCSD is that it is chemistry-guided rather than directly Hamiltonian-guided: the basis is restricted to the singles-and-doubles excitation structure expected to capture a large fraction of weak-to-moderate correlation effects.

The main parameters of UCCSD are the excitation amplitudes $\theta_i^a$ and $\theta_{ij}^{ab}$ attached to the singles-and-doubles pool. In a standalone VQE implementation, these amplitudes would be determined by variational energy minimization. In the present non-variational comparison, they are instead assigned classically using warm-start rules such as second-order M{\o}ller-Plesset perturbation theory (MP2) \cite{moller1934mp2}, which provides perturbative amplitudes, or coupled-cluster singles and doubles (CCSD) \cite{purvis1982ccsd}, which provides amplitudes from an iterative coupled-cluster calculation. Two UCCSD constructions are used in this work. The first is a one-shot full-versus-screened comparison, in which the full singles-and-doubles pool is evaluated directly against a screened or selected variant that retains only the larger-amplitude excitations. The second is the head-to-head $K$-dependent basis construction used later in the paper: starting from one MP2 or CCSD amplitude vector, the individual excitation generators are ranked by amplitude magnitude, and each added basis state is formed by applying one ranked generator at its fixed classical amplitude to the Hartree-Fock reference. This latter construction is not a conventional single-state UCCSD ansatz; rather, it is deliberately chosen to be closer in spirit to the GCIM idea \cite{zheng2024gcim}, where one builds a nonorthogonal basis from several UCC-generated states and then solves the projected eigenvalue problem in their span instead of optimizing one combined UCC circuit. Thus, MP2 or CCSD provides the amplitudes, while the full-versus-screened or generator-ranked choice determines how those amplitudes are turned into benchmark basis states.

\subsection{TrimCI-Based Reduced-Space Setup}

The working space is reduced to a TrimCI-selected determinant space before the Quantum Krylov and UCCSD benchmark states are constructed. The reason is that the expensive steps of the unreduced workflow scale with the size of the full many-body space. In particular, an explicit qubit-state representation of a generic $n$-qubit Quantum Krylov or UCCSD state scales as $2^n$, while exact FCI in the fixed-electron sector scales combinatorially with the number of spin orbitals and electrons, through the determinant-space dimension $\binom{n}{N_e}$. In the naive full-space setting of this project, this keeps the routine end-to-end benchmark only in the small-system regime, practically through LiH in the main workflow, with H$_2$O / STO-3G still manageable only as a standalone reference check. Such a regime is too small to support a meaningful extrapolation in qubit count, so larger systems are needed. In practice, this means that qubit-scaling trends should be probed at least beyond the 20-qubit scale, and preferably into the 40--50 qubit range. By projecting first into a compact selected determinant space, the main linear-algebra cost is controlled instead by the number of kept determinants rather than the full Hilbert-space dimension.

For a given molecular system, the first step is to run TrimCI in a shared restricted Hartree-Fock molecular-orbital basis and obtain an approximate ground state in determinant form,
\begin{equation}
\ket{\Psi_{\mathrm{TrimCI}}} = \sum_{\mu=1}^{M} c_{\mu} \ket{D_{\mu}}.
\end{equation}
In the present workflow, the practical TrimCI limiter is not an explicit energy-accuracy threshold but a determinant budget, controlled mainly by the parameter \texttt{ndets}, together with optional exploration controls such as \texttt{ndets\_explore} and \texttt{nexploration}. After the TrimCI run, the selected determinants are ranked by weight $|c_{\mu}|^2$.

The working subspace for the benchmark is then defined by truncating this list to the top $N$ determinants,
\begin{equation}
\mathcal{S}_{N} = \mathrm{span}\{\ket{D_{1}}, \ket{D_{2}}, \ldots, \ket{D_{N}}\},
\end{equation}
with $N \le M$. This additional cutoff removes the low-weight tail of the TrimCI expansion and gives a compact determinant basis for the later Quantum Krylov and UCCSD constructions. The adequacy of this reduced determinant space is validated explicitly against exact FCI on the smaller systems before the larger reduced-space benchmarks are interpreted.

All operators used in the benchmark are then represented directly in this $N$-determinant basis. The projected Hamiltonian is
\begin{equation}
H^{(N)}_{\mu \nu} = \bra{D_{\mu}} \hat{H} \ket{D_{\nu}},
\end{equation}
and the UCCSD excitation generators are represented in the same way,
\begin{equation}
G^{(a,N)}_{\mu \nu} = \bra{D_{\mu}} \left(\hat{\tau}_{a} - \hat{\tau}_{a}^{\dagger}\right) \ket{D_{\nu}}.
\end{equation}
Because the determinant basis itself is orthonormal, its intrinsic overlap matrix is the identity. The nontrivial overlap matrix appears only after the Quantum Krylov or UCCSD benchmark states are generated inside this reduced space.

Given these projected operators, the benchmark basis states are constructed entirely within $\mathcal{S}_{N}$. Quantum Krylov states are generated from the projected Hamiltonian and the projected Hartree-Fock reference. UCCSD states are generated from the projected excitation generators, either as one-shot full-versus-screened UCCSD states or, in the later head-to-head comparison, as ranked single-generator states built from MP2 or CCSD amplitudes. If a comparison uses $K$ benchmark basis states, then $K$ Quantum Krylov or UCCSD states are constructed in this reduced determinant space before the final projected eigensolver is applied.

\subsection{Benchmark Systems}

The benchmark ladder is chosen to serve two purposes at once: direct small-system validation against exact FCI and larger-system scaling tests in the TrimCI-reduced workflow. All systems are treated in the STO-3G basis. The smaller systems H$_2$, H$_4$, and LiH are used to calibrate and validate the reduced-space construction against exact reference data, while N$_2$, C$_2$H$_4$, and C$_3$H$_8$ extend the study into a regime where full-space simulation would already be impractical in the end-to-end workflow but the TrimCI-based reduced-space approach remains tractable.

\begin{table}[t]
\caption{Benchmark systems used in the reduced-space study. All systems are treated in the STO-3G basis. The Hamiltonian Pauli-term column reports the number of terms in the mapped qubit Hamiltonian; the C$_3$H$_8$ entry is listed as a coarse extrapolated estimate because direct compilation of the full 46-qubit Hamiltonian was prohibitively slow in the present environment.}
\label{tab:benchmark-systems}
\begin{ruledtabular}
\scriptsize
\setlength{\tabcolsep}{3pt}
\begin{tabular}{@{}ccccc@{}}
System & Qubits & Electrons & \makecell[c]{Hamiltonian\\Pauli terms} & Role \\
\midrule
H$_2$ & 4 & 2 & 15 & \makecell[c]{Exact-FCI\\calibration} \\
H$_4$ & 8 & 4 & 185 & \makecell[c]{Exact-FCI\\calibration} \\
LiH & 12 & 4 & 631 & \makecell[c]{Exact-FCI\\calibration} \\
N$_2$ & 20 & 14 & 3,327 & \makecell[c]{Reduced-space\\scaling} \\
C$_2$H$_4$ & 28 & 16 & 8,919 & \makecell[c]{Reduced-space\\scaling} \\
C$_3$H$_8$ & 46 & 26 & $\sim 3.8\times 10^4$ & \makecell[c]{Reduced-space\\scaling} \\
\end{tabular}
\end{ruledtabular}
\end{table}

\subsection{Projected Subspace Construction}

All basis families start from the same Hartree-Fock reference state $\ket{\phi_{\mathrm{HF}}}$. Given a set of prepared basis states $\{\ket{\phi_i}\}_{i=0}^{K-1}$, the projected Hamiltonian and overlap matrices are formed as
\begin{align}
H_{ij} &= \bra{\phi_i}\hat{H}\ket{\phi_j}, \\
S_{ij} &= \braket{\phi_i|\phi_j},
\end{align}
and the lowest generalized eigenvalue of
\begin{equation}
Hc = E_{\mathrm{sub}} Sc
\end{equation}
is used as the subspace energy estimate.

\subsection{Generalized Eigensolver and Overlap Conditioning}

The projected problem takes the generalized form rather than the ordinary form because the Quantum Krylov and UCCSD basis states are generally nonorthogonal. Thus the subspace coefficients satisfy $Hc = E_{\mathrm{sub}}Sc$, and this reduces to an ordinary eigenvalue problem only in the special case of an orthonormal basis with $S = I$.

Because the nonorthogonal basis can contain nearly dependent directions, the overlap matrix may become ill-conditioned. To control this, $S$ is eigendecomposed as
\begin{equation}
S = U \, \mathrm{diag}(s_i) \, U^{\dagger},
\end{equation}
and directions with overlap eigenvalues below a threshold are discarded. In the current implementation, this is done with an absolute eigenvalue test $s_i > s_{\mathrm{eval\_cutoff}}$ rather than with a relative rule. The released benchmark notebooks use $s_{\mathrm{eval\_cutoff}} = 10^{-8}$ in one notebook and a looser $10^{-10}$ cutoff in the H$_4$ and LiH notebooks, so the working basis is orthogonalized only on the numerically stable part of the span. The transformation
\begin{equation}
X = U_{\mathrm{keep}} \, \mathrm{diag}(s_i^{-1/2})
\end{equation}
then maps the retained nonorthogonal basis to an orthonormal basis, and the Hamiltonian is transformed as
\begin{equation}
\widetilde{H} = X^{\dagger} H X .
\end{equation}
The final projected estimate is obtained by solving the ordinary eigenvalue problem for $\widetilde{H}$ in this retained orthogonal subspace.

In a hardware-oriented execution, the matrix elements of $H$ and $S$ would have to be obtained from Hadamard tests or related overlap-estimation techniques. In the present work, however, the purpose is to compare Quantum Krylov and UCCSD as basis generators rather than to benchmark matrix-element measurement protocols. The benchmark states are therefore represented analytically in the reduced TrimCI space, and the entries of $H$ and $S$ are formed directly from inner products and projected operators before solving the generalized eigenvalue problem. For larger reduced-space instances, iterative eigensolvers such as Davidson or LOBPCG are natural choices.

\subsection{Comparison Protocol}

The comparison in this work is non-variational for both basis families. Quantum Krylov and UCCSD are not judged by the result of a full parameter-optimization loop, but by the projected subspaces generated from their method-specific state-construction rules. This choice isolates the basis generator itself rather than the trainability of a variational ansatz.

Under this protocol, both methods are evaluated for the same molecular Hamiltonian, in the same orbital basis, from the same Hartree-Fock reference state, inside the same TrimCI-defined reduced determinant space, and at the same target benchmark-basis size $K$. The projected ground-state estimate is always obtained from the same Rayleigh-Ritz procedure together with the same overlap-conditioning procedure. Exact FCI is used as the validation reference on the smaller calibration systems, while the larger systems are compared within the same TrimCI-reduced workflow validated against those smaller exact benchmarks.

The two families are allowed to keep their natural parameterization. For Quantum Krylov, the defining parameter is the time spacing $\Delta t$ that generates the equally spaced evolution times $t_k = k\Delta t$ used for the time-evolved basis states. In this paper, $\Delta t$ is selected by grid search. For UCCSD, the defining ingredients are the amplitudes attached to the singles-and-doubles excitation pool together with the choice of whether to keep the full pool or a screened subset. In the present non-variational setting, these amplitudes are supplied by classical chemistry information such as MP2 or CCSD rather than by a VQE optimization loop. The comparison therefore does not claim to identify the best standalone Quantum Krylov or the best standalone UCCSD-VQE result. Instead, it asks which family provides the more expressive projected subspace, within the same reduced-space representation, for the cost of preparing its states and operators.

For the explicit $K$-by-$K$ head-to-head comparison reported later, the UCCSD side is organized as two separate ranked generator families, one built from MP2 amplitudes and one built from CCSD amplitudes. In each family, the individual singles-and-doubles generators are ranked by amplitude magnitude, and each added basis state is formed by applying one ranked generator at its fixed classical amplitude to the Hartree-Fock reference. Thus $K=1$ is the Hartree-Fock state itself, $K=2$ adds the top-ranked single-generator state, $K=3$ adds the second-ranked single-generator state, and so on. This makes the later UCCSD comparison closer in spirit to a GCIM-style \cite{zheng2024gcim} nonorthogonal basis built from UCC generators than to a nested full-UCCSD truncation ladder. Curves may terminate before the target $K$ if the overlap-conditioning rule rejects additional nearly dependent states.

\subsection{Assumptions and Parameters}

Several assumptions are held fixed so that the ansatz comparison remains controlled.

\begin{itemize}
    \item The main benchmark is noiseless and focuses on basis quality and preparation cost, not hardware noise or full shot-noise modeling.
    \item The primary fairness axis is basis size $K$; total trainable parameter count is tracked only as a secondary diagnostic.
    \item Quantum Krylov and UCCSD use the same molecular Hamiltonian, the same Hartree-Fock reference state, the same projected-subspace objective, and the same validation reference.
    \item The comparison is non-variational. Quantum Krylov basis states are specified by chosen evolution times, while UCCSD basis states are specified either as one-shot full-versus-screened ans\"atze or as ranked single-generator states built from MP2 or CCSD amplitudes.
    \item Candidate states are rejected when the overlap matrix becomes numerically ill-conditioned. In the current code, this is implemented with an absolute smallest-eigenvalue cutoff on $S$, with $s_{\mathrm{eval\_cutoff}}=10^{-8}$ in one benchmark notebook and $10^{-10}$ in the looser H$_4$ and LiH notebooks.
    \item Small calibration systems are used intentionally so that the reduced-space workflow can be checked against exact FCI data before moving to the larger systems.
\end{itemize}

\section{Results and Validation}

This section is organized around the method-specific choices that remain in the non-variational comparison. On the Krylov side, the main task is to select a useful time spacing by grid search. On the UCCSD side, the main task is to decide whether the full singles-and-doubles pool is necessary, or whether a screened subset of important excitations is already sufficient when amplitudes are assigned classically by MP2 or CCSD. After those two ingredients are fixed, the final subsection compares the resulting Krylov and UCCSD basis families directly as the benchmark basis size $K$ increases. The supporting validation that the TrimCI reduced-space reference reproduces exact FCI on the smaller calibration systems is provided in Appendix~\ref{subsec:trimci-validation}.

\subsection{Krylov Time-Step Selection by Grid Search}

Even though the benchmark avoids a full variational optimization loop, the Krylov basis still depends on a method-specific parameter, namely the time spacing $\Delta t$ used to define the equally spaced sequence $t_k = k\Delta t$. For this reason, the first stage of the results is a grid-search comparison over candidate $\Delta t$ values. For each molecule, candidate spacings are scanned, choices that lead to an ill-conditioned overlap matrix are discarded, and the remaining spacing with the lowest projected ground-state energy is retained. In the reduced-space workflow used here, the target Krylov size is capped at $K \le 4$; for H$_2$, the retained TrimCI space is only two-dimensional, so the largest nontrivial Krylov basis is $K=2$.

\begin{figure}[t]
    \includegraphics[width=\columnwidth]{../P765_final_project/Krylov_Grid_Search_N2.pdf}
    \caption{Representative Krylov grid search for N$_2$ in the TrimCI-reduced space. The plotted quantity is the projected energy gap $E_{\mathrm{Krylov}}-E_{\mathrm{ref}}$ as a function of the spacing $\Delta t$ at $K=4$. Red crosses mark candidate spacings rejected by the overlap cutoff, and the vertical dotted line marks the selected spacing.}
    \label{fig:krylov-grid-search-n2}
\end{figure}

Figure~\ref{fig:krylov-grid-search-n2} shows the representative N$_2$ scan. The smallest spacings are rejected because the overlap matrix becomes too ill-conditioned, while a shallow minimum appears once the time-evolved states become sufficiently distinct. The best spacing is $\Delta t = 0.150$ for N$_2$, with a reduced-space Krylov gap of $0.698$ mHa relative to the TrimCI reference. The same grid-search protocol was applied to every benchmark system, and the resulting best spacings are summarized in the table below. A striking feature of the summary is that the optimal spacing is very stable across the ladder: $\Delta t = 0.150$ is selected for H$_4$, LiH, N$_2$, and C$_2$H$_4$, while the largest system C$_3$H$_8$ shifts only slightly to $\Delta t = 0.300$.

\begin{center}
    \small
    \begin{tabular}{lcccc}
        \toprule
        System & Qubits & $N_{\mathrm{det}}$ & Best $\Delta t$ & Best gap (mHa) \\
        \midrule
        H$_2$ & 4 & 2 & 0.100 & -0.000 \\
        H$_4$ & 8 & 20 & 0.150 & 0.244 \\
        LiH & 12 & 69 & 0.150 & 0.631 \\
        N$_2$ & 20 & 73 & 0.150 & 0.698 \\
        C$_2$H$_4$ & 28 & 73 & 0.150 & 0.584 \\
        C$_3$H$_8$ & 46 & 71 & 0.300 & 0.408 \\
        \bottomrule
    \end{tabular}
\end{center}

\subsection{Full and Screened UCCSD with MP2/CCSD Parameters}

The second stage of the results examines whether the full UCCSD excitation pool is actually needed once amplitudes are fixed by classical chemistry information. To do this, I evaluate non-variational UCCSD states at MP2 and CCSD amplitudes, and then compare the full singles-and-doubles pool against a screened variant that keeps only operators with amplitude magnitude above the fixed threshold $10^{-2}$. In this work, the screened ansatz is the selected-UCC surrogate used for hardware-cost reduction. Adaptive UCC variants are not included, because they would reintroduce an iterative variational loop outside the intended scope of this non-variational comparison. The detailed screening study is reported on H$_2$, H$_4$, and LiH, where exact FCI remains available and full one-shot UCCSD evaluation is still routine in the present implementation.

\begin{figure}[t]
    \includegraphics[width=\columnwidth]{../P765_final_project/UCCSD_Screening_LiH_Combined.pdf}
    \caption{Representative LiH comparison for full and screened one-shot no-opt UCCSD with MP2- and CCSD-parameterized amplitudes. Top: energy gaps relative to exact FCI, with the dashed line marking chemical accuracy. Bottom: retained excitation-operator counts for the same four variants.}
    \label{fig:uccsd-screening-lih}
\end{figure}

Figure~\ref{fig:uccsd-screening-lih} shows the representative LiH tradeoff. In the full UCCSD pool, the MP2- and CCSD-parameterized states have FCI gaps of $11.224$ and $9.002$ mHa, respectively. After screening, those gaps change only slightly, to $11.308$ mHa for MP2 and $9.073$ mHa for CCSD, while the operator counts collapse from $92$ to $6$ and $8$. This is the central screening result of the subsection: at least on the calibration systems, most of the one-shot UCCSD benefit survives after discarding the small-amplitude tail of the excitation pool.

The same pattern appears across the rest of the small-system set. For H$_2$, screening reduces the operator count from $3$ to $1$ with no visible change in the reported gap. For H$_4$, screening reduces the operator count from $26$ to $10$ while changing the gap by less than $0.03$ mHa. In all three systems, the CCSD-initialized ansatz is more accurate than the MP2-initialized ansatz, but the main hardware-relevant message is that the screened construction tracks the full construction very closely at substantially lower operator count.

\begin{table}[t]
    \centering
    \small
    \begin{tabular}{llcc}
    \toprule
    System & Variant & Gap to FCI (mHa) & Ops \\
    \midrule
    H$_2$ & MP2 full & 9.473 & 3 \\
     & MP2 screened & 9.473 & 1 \\
     & CCSD full & 5.104 & 3 \\
     & CCSD screened & 5.104 & 1 \\
    H$_4$ & MP2 full & 23.613 & 26 \\
     & MP2 screened & 23.613 & 10 \\
     & CCSD full & 18.793 & 26 \\
     & CCSD screened & 18.768 & 10 \\
    LiH & MP2 full & 11.224 & 92 \\
     & MP2 screened & 11.308 & 6 \\
     & CCSD full & 9.002 & 92 \\
     & CCSD screened & 9.073 & 8 \\
    \bottomrule
\end{tabular}

    \caption{Full and screened one-shot UCCSD variants on the small calibration systems. Reported gaps are relative to exact FCI, and `Ops' counts the retained excitation operators.}
    \label{tab:uccsd-screening-summary}
\end{table}

For the later cross-family comparison, this subsection therefore motivates a nested UCCSD construction rather than a single fixed UCCSD state. In that later comparison, the excitation amplitudes are ranked by importance and used to build growing MP2-based and CCSD-based UCCSD hierarchies, so that the UCCSD side can also be studied as a function of basis size $K$.

\subsection{Head-to-Head Krylov versus UCCSD as $K$ Increases}

With the Krylov spacing fixed by grid search, the final stage is a direct comparison of how the basis families themselves improve as $K$ grows across the full six-system benchmark ladder. The figure therefore spans H$_2$, H$_4$, LiH, N$_2$, C$_2$H$_4$, and C$_3$H$_8$, so that the head-to-head comparison can be read from the smallest to the largest reduced-space instances used in this work. In every case, the projected energy is reported relative to the same TrimCI reduced-space reference energy. On the Krylov side, the plotted curve uses the best grid-searched spacing for the largest numerically stable Krylov target up to $K=20$. On the UCCSD side, I use two separate non-variational generator-ranked families, one ranked by MP2 amplitudes and one ranked by CCSD amplitudes. Each added UCCSD basis state is generated from a single ranked excitation operator acting on the Hartree-Fock reference at its fixed classical amplitude, so the ordered sequence is HF at $K=1$, then the top-ranked single-generator state, then the second-ranked single-generator state, and so on. In this sense, the head-to-head UCCSD comparison is closer to a GCIM-style nonorthogonal basis built from UCC generators than to a nested truncation of one full UCCSD circuit \cite{zheng2024gcim}.

\begin{figure*}[t]
    \centering
    \includegraphics[width=\textwidth]{../P765_final_project/Krylov_vs_UCCSD_K_Comparison.pdf}
    \caption{Head-to-head comparison of Krylov, MP2, and CCSD basis families as the benchmark basis size $K$ increases across the full benchmark ladder. The vertical axis shows the absolute projected-energy gap $|E_{\mathrm{sub}}-E_{\mathrm{ref}}|$ on a logarithmic scale, with values below $10^{-6}$ mHa clipped to that floor for display. The gray shaded region marks gaps below chemical accuracy ($1.6$ mHa). `MP2' and `CCSD' denote the generator-ranked UCC basis families built from fixed MP2 or CCSD amplitudes, respectively, where $K=1$ is Hartree-Fock and each additional basis state is generated by one more ranked single UCC excitation generator. The horizontal dotted lines use the same MP2 and CCSD colors and mark the corresponding full-pool one-shot UCCSD baselines in the same reduced determinant space. Curves terminate when the overlap-conditioning rule rejects further nearly dependent states or when the available reduced-space dimension is exhausted.}
    \label{fig:krylov-vs-uccsd-k}
\end{figure*}

Figure~\ref{fig:krylov-vs-uccsd-k} shows that the three families behave quite differently once the comparison is phrased in terms of how quickly their projected span improves with $K$. Krylov is the only family that consistently drives the projected energy close to the reduced-space reference across the full ladder: it is exact for H$_2$ by $K=2$, reaches the numerical floor for H$_4$ by $K=12$, and then continues to converge to $10^{-4}$--$10^{-6}$ mHa levels on LiH, N$_2$, C$_2$H$_4$, and C$_3$H$_8$. The best Krylov spacings selected by the grid search are $\Delta t^*=0.10$ (H$_2$), $1.95$ (H$_4$), $1.70$ (LiH), $1.05$ (N$_2$), $1.20$ (C$_2$H$_4$), and $0.95$ (C$_3$H$_8$), while the largest numerically stable Krylov targets range from $K=2$ for H$_2$ up to the full $K=20$ target on N$_2$, C$_2$H$_4$, and C$_3$H$_8$.

The UCC side shows a different pattern. On the smallest systems the generator-ranked MP2 and CCSD families can also approach the reference, but as the systems grow they converge much more slowly and often terminate early because of overlap conditioning. The CCSD-ranked family is usually the stronger of the two on the small calibration systems, reaching $0.287$ mHa on H$_4$ and $0.0209$ mHa on LiH, but it still remains at $20.09$ mHa on N$_2$, $21.16$ mHa on C$_2$H$_4$, and $33.47$ mHa on C$_3$H$_8$ at the largest sampled $K$. The MP2-ranked family stays above chemical accuracy as well on the larger systems, with final gaps of $19.82$ mHa on N$_2$, $11.92$ mHa on C$_2$H$_4$, and $21.98$ mHa on C$_3$H$_8$. The main conclusion from the six-system figure is therefore that Krylov gains much more from increasing $K$ than the amplitude-ranked single-generator UCC families once the problem moves beyond the small calibration regime. At the same time, the spread between the MP2 and CCSD curves shows that, in this GCIM-like UCC construction, the parent classical amplitude source and the conditioning of the resulting nonorthogonal basis are both essential parts of the method design.

\subsection{Extrapolated Basis Size at Chemical Accuracy}

The qubit-scaling trend can be summarized more compactly by extracting, from Fig.~\ref{fig:krylov-vs-uccsd-k}, the basis size required to reach chemical accuracy. For each system and basis family, I define $K_{\mathrm{chem}}$ as the first sampled $K$ for which $|E_{\mathrm{sub}}-E_{\mathrm{ref}}|\le 1.6$ mHa. If the sampled data do not reach that threshold, I fit the tail of the convergence curve to the log-linear form
\begin{equation}
\log_{10}|E_{\mathrm{sub}}-E_{\mathrm{ref}}| = a + bK,
\end{equation}
using the last up to five available points, and solve for the crossing with $1.6$ mHa to obtain an extrapolated estimate of $K_{\mathrm{chem}}$.

\begin{figure}[t]
    \centering
    \includegraphics[width=\columnwidth]{../P765_final_project/K_Chemical_Accuracy_Extrapolation.pdf}
    \caption{Estimated basis size required to reach chemical accuracy as a function of qubit count. Filled markers denote systems where the sampled data already reach the $1.6$ mHa target, while open markers denote extrapolated values obtained from a log-linear fit to the tail of the convergence curve.}
    \label{fig:k-chemical-accuracy}
\end{figure}

Figure~\ref{fig:k-chemical-accuracy} makes the scaling separation between the basis families explicit. Krylov reaches chemical accuracy with only a handful of basis states across the full ladder: $K_{\mathrm{chem}}=2$ for H$_2$, $11$ for H$_4$, $2$ for LiH, $5$ for N$_2$, and $4$ for both C$_2$H$_4$ and C$_3$H$_8$. The generator-ranked UCC families behave very differently. On the small systems, the MP2 and CCSD constructions still reach chemical accuracy with moderate basis sizes, namely $(2,11,5)$ for MP2 and $(2,10,7)$ for CCSD on H$_2$, H$_4$, and LiH. Beyond that regime, however, chemical accuracy is not reached within the sampled $K$ window, and the extrapolated requirements rise into the tens of basis states. The extrapolated values are approximately $K_{\mathrm{chem}}\approx 44$, $34$, and $63$ for the MP2 family on N$_2$, C$_2$H$_4$, and C$_3$H$_8$, and $K_{\mathrm{chem}}\approx 47$, $41$, and $85$ for the corresponding CCSD family. This final scaling summary is consistent with the direct convergence plots: within the present reduced-space benchmark, Krylov maintains a much more favorable expressivity-per-basis trend than the non-variational UCC constructions as the system size increases.

\section{Full-Stack Impact / Resource Implications}

This section is required for the final project. State clearly how your result or tool affects a stack-level conclusion about utility-scale quantum computing under explicit assumptions.

\subsection{Resource Estimation}

Because the dependence on the number of basis states $K$ has already been studied in Sec.~IV, the resource comparison here focuses on the intrinsic cost of preparing one non-Hartree-Fock basis state. For Krylov, the reported values are based on first-order Trotterized time evolution with $n=2$ steps at the benchmark-selected $\Delta t^*$. For UCCSD-MP2 and UCCSD-CCSD, the reported values are averaged over the ranked single-generator basis states used in the head-to-head comparison. The combined table below summarizes total gate count, circuit depth, and two-qubit-gate count per basis state. The C$_3$H$_8$ Krylov row is an extrapolated proxy based on the exact smaller-system trend, while the remaining rows are taken directly from the current PennyLane resource estimates.

\begin{table*}[t]
    \centering
    \scriptsize
    \resizebox{\textwidth}{!}{\begin{tabular}{lrrrrrrrrrrrrrrrrrr}
    \toprule
    & \multicolumn{3}{c}{H$_2$} & \multicolumn{3}{c}{H$_4$} & \multicolumn{3}{c}{LiH} & \multicolumn{3}{c}{N$_2$} & \multicolumn{3}{c}{C$_2$H$_4$} & \multicolumn{3}{c}{C$_3$H$_8$} \\
    \cmidrule(lr){2-4} \cmidrule(lr){5-7} \cmidrule(lr){8-10} \cmidrule(lr){11-13} \cmidrule(lr){14-16} \cmidrule(lr){17-19}
    Method & T & Depth & 2Q & T & Depth & 2Q & T & Depth & 2Q & T & Depth & 2Q & T & Depth & 2Q & T & Depth & 2Q \\
    \midrule
    Krylov & 1232 & 26 & 12 & 16192 & 362 & 56 & 55440 & 1250 & 148 & 261536 & 5926 & 380 & 784784 & 17774 & 756 & 4235201.0 & 98801.3 & 2289.4 \\
    UCCSD-MP2 & 352 & 20 & 14 & 352 & 20 & 14 & 352 & 20 & 14 & 352 & 20 & 14 & 352 & 20 & 14 & 352 & 20 & 14 \\
    UCCSD-CCSD & 352 & 20 & 14 & 277.7 & 15.7 & 10.6 & 297.3 & 16.8 & 11.5 & 352 & 20 & 14 & 352 & 20 & 14 & 352 & 20 & 14 \\
    \bottomrule
\end{tabular}
}
    \caption{Per-basis-state resource summary for the current Krylov, UCCSD-MP2, and UCCSD-CCSD constructions. Each method block reports the average total native gate count, circuit depth, and two-qubit-gate count for one non-Hartree-Fock basis state. The C$_3$H$_8$ Krylov entry is an extrapolated proxy based on the exact smaller-system trend; all remaining entries are drawn directly from the current PennyLane resource estimates.}
    \label{tab:resource-combined}
\end{table*}

\begin{figure}[t]
    \centering
    \includegraphics[width=\columnwidth]{../P765_final_project/Resource_PerState_DualAxis.pdf}
    \caption{Per-basis-state resource scaling versus qubit count. Colors distinguish Krylov, UCCSD-MP2, and UCCSD-CCSD. On the left axis, solid-circle lines show the total native gate count and dotted-square lines show the two-qubit-gate count. On the right axis, dashed-triangle lines show the circuit depth. The C$_3$H$_8$ Krylov point is an extrapolated proxy, while the remaining values are taken directly from the current PennyLane resource estimates.}
    \label{fig:resource-dual-axis}
\end{figure}

For LiH, the same resource picture is reflected in the comparison between full
and amplitude-selected UCCSD. The selected construction preserves the one-shot
warm-start energy to within about $0.1$~mHa while reducing the compiled circuit
cost by more than an order of magnitude. In the MP2 branch, the selected
construction lowers the gate count from $9344$ to $720$, and in the CCSD branch
from $9344$ to $748$, while the energy gap changes only from $11.22$ to
$11.31$~mHa and from $9.00$ to $9.07$~mHa, respectively. This makes selected
UCCSD the clearly preferable chemistry-inspired basis when the objective is to
reduce circuit cost without reintroducing a variational loop.

The corresponding selected-Krylov truncation study is moved to
Appendix~\ref{subsec:selected-krylov-appendix}, where it is kept as a supporting
resource tradeoff rather than a main-text result.

\subsection{Full-Stack Impact}

The stack-level interpretation is more nuanced than a single "winner" between
Krylov and UCCSD. At the quantum-algorithm layer, Krylov remains attractive
because its stronger expressivity lowers the basis size required for chemical
accuracy, as shown in Fig.~\ref{fig:k-chemical-accuracy}. Since the number of
projected Hamiltonian and overlap elements grows roughly as $K^2$, this also
reduces the shot burden of the subspace solve. Using the extracted
$K_{\mathrm{chem}}$ values, the number of unique basis pairs
$K_{\mathrm{chem}}(K_{\mathrm{chem}}+1)/2$ is only $15$ for Krylov on N$_2$ and
$10$ for Krylov on both C$_2$H$_4$ and C$_3$H$_8$, compared with $990$ and
$2016$ for the MP2-ranked UCC family and $1128$ and $3655$ for the
CCSD-ranked UCC family on N$_2$ and C$_3$H$_8$, respectively. In this sense,
Krylov is algorithmically more efficient: it reaches the same reduced-space
accuracy with a much smaller projected subspace.

At the logical-processor layer, however, Fig.~\ref{fig:resource-dual-axis}
shows the opposite trend: each Krylov basis state is much deeper than the
corresponding selected UCC basis state because it requires Trotterized
Hamiltonian evolution. A simple depth-dominated logical-fidelity estimate can
be obtained by requiring the total logical failure probability for preparing the
$K_{\mathrm{chem}}$ basis states to remain below a budget $\eta$,
\begin{equation}
    p_L^{\mathrm{req}} \approx \frac{\eta}{K_{\mathrm{chem}} d_{\mathrm{state}}},
\end{equation}
where $d_{\mathrm{state}}$ is the per-state circuit depth. In
this estimate, I set $\eta=10^{-2}$ as a rough one-pass budget for basis-state
preparation. The estimate is optimistic because it ignores the repeated
sampling needed to measure the projected $H$ and $S$ matrices, but it is still
sufficient to expose the stack-level trend.

This estimate shows that selected UCCSD is substantially easier at the logical
and error-correction layers. On N$_2$, for example, the resulting requirement is
about $3.0\times 10^{-7}$ for Krylov, compared with $1.1\times 10^{-5}$ for
both MP2- and CCSD-ranked UCC. On C$_3$H$_8$, the corresponding values are
$2.4\times 10^{-8}$, $7.9\times 10^{-6}$, and $5.9\times 10^{-6}$. Under any
standard fault-tolerant architecture, a smaller target logical error rate
translates into a larger code distance and therefore a heavier physical-processor
overhead. In that sense, selected UCCSD is favored at the logical and physical
layers even though it is less expressive at the algorithmic layer.

The overall full-stack conclusion is therefore mixed rather than one-sided.
Krylov suppresses the basis-size and shot-scaling burden, while selected UCCSD
relaxes the required logical fidelity per circuit. Which family is closer to
practical implementation depends on the dominant bottleneck: if projected-subspace
size and sampling cost dominate, Krylov is more appealing; if logical depth and
error-correction overhead dominate, selected UCCSD is more realistic. This
benchmark therefore suggests that the ansatz choice is not just an algorithmic
question, but a stack-wide co-design question linking quantum chemistry, logical
resources, and hardware feasibility.

\section{Conclusion}

This work compared Quantum Krylov and UCCSD as non-variational basis generators for
molecular ground-state approximation in a projected-subspace framework. By
placing both families in the same reduced determinant-space workflow, the
benchmark focused on basis quality rather than on a variational optimization
loop. The main numerical result is that Quantum Krylov reaches chemical accuracy with a
substantially smaller basis size than the non-variational UCC constructions
across the benchmark ladder, indicating a stronger expressivity-per-basis trend
as the system size grows.

At the same time, the resource analysis shows that this algorithmic advantage
does not translate into a simple hardware advantage. Quantum Krylov basis states are
much more expensive to prepare because they require Trotterized Hamiltonian
evolution, whereas selected UCCSD states remain shallow and comparatively
cheap. The resulting picture is therefore a genuine tradeoff: Quantum Krylov is favored
when the dominant concern is projected-subspace size and sampling burden, while
selected UCCSD is favored when the dominant concern is logical depth and
hardware cost per basis state.

Taken together, these results suggest that basis design for quantum chemistry
should be evaluated at more than one layer of the stack. In the present
benchmark, Quantum Krylov is the more powerful basis family at the algorithmic level,
but selected UCCSD remains attractive from the perspective of implementation.
This makes the choice between them not just a question of accuracy, but a
co-design problem linking ansatz expressivity, logical resources, and hardware
feasibility.

\section{Limitations and Future Work}

\subsection{Current Limitations}

The main limitation of the present benchmark is that the working subspace is
not the full FCI space for the entire system ladder. Instead, the Hamiltonian,
reference energy, and benchmark basis states are all defined in a
TrimCI-selected determinant space. The reduced space reproduces the exact FCI
result very well on the small calibration systems, but the larger-system
conclusions in this report should still be interpreted as reduced-space
results rather than exact full-space statements.

A second limitation is that the resource comparison is not yet fully symmetric
between the two basis families. In the current implementation, Quantum Krylov basis
states are generated from Trotterized Hamiltonian evolution, whereas the UCCSD
side is represented by selected single-generator states ranked by MP2 or CCSD
amplitudes. This is a meaningful practical comparison, but it does not exhaust
all possible hardware-aware variants of either family. In particular, a
truncated or selected Quantum Krylov construction could make the resource comparison
more balanced against selected UCCSD.

\subsection{Future Work}

One natural extension is to make the hardware-level benchmark more symmetric.
This includes repeating the resource analysis with truncated-Hamiltonian Quantum Krylov
constructions and comparing those results directly against warm-start and
selected-UCCSD variants. Such a study would clarify whether lower-cost Quantum Krylov
states can preserve enough of the expressivity advantage seen in the present
benchmark. A closely related direction is to investigate whether qubitization-
based implementations of Quantum Krylov time evolution can reduce the resource cost of
the Quantum Krylov basis relative to the current Trotterized construction.

A second direction is to expand the choice of reference state beyond
Hartree-Fock. For example, both Quantum Krylov and UCCSD could be initialized from a
stronger warm-start state such as the TrimCI reduced-space ground state rather
than the Hartree-Fock determinant. That would help separate the effect of basis
generation from the effect of the initial reference itself.

It would also be valuable to connect the present resource estimates more
directly to concrete hardware specifications. Mapping the benchmark circuits to
specific platform assumptions, such as superconducting or trapped-ion devices,
would make it possible to assess which basis family is closer to realistic
implementation on a given hardware stack.

Finally, the present comparison suggests several algorithmic extensions. One
possibility is to design a new hybrid basis family that combines the compact
expressivity of Quantum Krylov with the shallower state-preparation structure of
selected UCCSD. Another is to investigate related subspace constructions and
variations of UCC/Quantum Krylov ideas, for example Q-SENSE
\cite{patel2025qsense}. A complementary theoretical direction is to study more
carefully why Quantum Krylov bases so effectively describe the ground state in these
electronic-structure problems, and whether that behavior can be predicted from
spectral or correlation properties of the Hamiltonian.

\section{Repository and Reproducibility}

This project is organized around a small set of benchmark notebooks and validation scripts. The main software stack uses PennyLane for quantum-state construction, PySCF for classical electronic-structure reference calculations, and TrimCI for reduced determinant-space selection.

\textbf{Repository:} \url{https://github.com/jhp1859/p759_final_project/draft_deliveries}

\textbf{Suggested reproduction steps:}
\begin{enumerate}
    \item Clone or download the repository.
    \item Create the project environment and install the required packages for PennyLane, PySCF, and TrimCI.
    \item Run \texttt{validate\_trimci\_subspace.py} to reproduce the TrimCI-versus-FCI reference-validation figure and CSV file.
    \item Run the benchmark notebooks for H$_2$, H$_4$, and LiH to reproduce the small-system Krylov and UCCSD comparisons.
    \item Run the larger-system reduced-space scripts to reproduce the warm-start and H$_2$O comparison figures used to motivate the scaling discussion.
\end{enumerate}

\textbf{Main entry points:}
\begin{itemize}
    \item \texttt{Fair\_Subspace\_Benchmark.ipynb}, \texttt{Fair\_Subspace\_Benchmark\_H4.ipynb}, and \texttt{Fair\_Subspace\_Benchmark\_LiH.ipynb}: end-to-end small-system benchmarks for Krylov and UCCSD.
    \item \texttt{validate\_trimci\_subspace.py}: validation of the TrimCI reduced-space reference against exact FCI, together with the figure \texttt{TrimCI\_FCI\_Subspace\_Validation.pdf}.
    \item \texttt{warmstart\_uccsd\_compare.py}: MP2- and CCSD-parameterized full and screened UCCSD warm-start comparisons.
    \item \texttt{uccsd\_screening\_benchmark.py}: compact full-versus-screened UCCSD benchmark producing the LiH representative figure and the H$_2$/H$_4$/LiH summary table.
    \item \texttt{krylov\_vs\_uccsd\_k\_comparison.py}: three-system head-to-head comparison of Krylov and nested MP2/CCSD-ranked UCCSD hierarchies as the basis size $K$ increases.
    \item \texttt{h2o\_trimci\_krylov\_uccsd\_compare.py}: larger-system H$_2$O / STO-3G comparison in the TrimCI-reduced workflow.
    \item \texttt{Resource\_Estimation\_Subspace.ipynb}: resource-estimation summaries for the benchmark ansatz families.
\end{itemize}

% Uncomment and fill in if you worked in a pair.
% \section*{Contribution Statement}
% [Briefly state who implemented what.]


\clearpage
\bibliography{references}

\end{document}
%
% ****** End of file final_project_report.tex ******
